\section{Conclusion}
\label{sec:conclusion}

This paper presents a defense-oriented benchmark for evaluating security systems against multi-stage APT campaigns. Rather than treating attack techniques or alerts as independent evaluation samples, the framework preserves individual playbook action occurrences, records the defense outcome of each action, and binds that observation to a normalized scoring opportunity representing the first failed victim-side control point. Canonical attack phases, defense-opportunity mappings, and stage dependency semantics then connect local observations to stage-, playbook-, and benchmark-level scores.

The resulting design makes two aspects of APT defense evaluation explicit. First, defensive effectiveness must be interpreted in the context of both attack-chain structure and the control surface of the evaluated system; raw detection counts alone cannot express whether a viable attack path remains. Second, an aggregate score is useful only when it remains traceable to the evidence that produced it. Our benchmark therefore retains a reverse path from an overall result to the corresponding playbook, stage, raw action, observed defense response, scoring opportunity, and repairable control failure.

We construct the benchmark around 53 APT playbooks and preserve all 1,796 authoritative raw action occurrences in the experimental protocol. Our ongoing controlled evaluation uses real defense-system observations and explicitly records replay fidelity for actions that require emulation or controlled substitution. [FINAL RESULT SENTENCE(S): insert the principal empirical comparison and the strongest repair-oriented finding after completion of the real experiments.]

More broadly, this work argues that APT defense evaluation should move beyond asking only whether an attack was detected. A useful benchmark should also determine where the attack chain remained executable, why the defense failed at that point, and which control should be strengthened to prevent the same progression. By preserving this connection between attack execution, defense observation, quantitative scoring, and remediation, the proposed framework aims to make APT defense evaluation both comparable and operationally actionable.